\section{Introduction}
\label{sec:intro}

Mapping malware evidence to MITRE ATT\&CK technique identifiers is a labelling task that large
language models are now routinely applied to, and the systems that perform it are rarely a model
alone. A working pipeline wraps the model in deterministic tooling: a sandbox that detonates the
sample, a disassembler that recovers its structure, retrieval indices that supply prior cases, and
rule engines that assert techniques directly from signatures. The model's contribution is therefore
a difference against that tooling rather than a score in isolation, and that difference is almost
never reported. Büchel et al.'s systematisation re-evaluates more than forty technique-extraction
systems in one setting and finds them ``largely incomparable'', because each uses its own ontology
and its own inaccessible dataset \cite{buchel2025sok}. An F1 value read against no baseline refers
to nothing, and the baseline that matters for an LLM pipeline is the deterministic engine it was
built on top of.

A second barrier is that such a pipeline is itself a distributed system, and its evaluation inherits
every property of one. Each boundary between the pipeline and a server it depends on is a place
where a request can succeed and still not mean what the caller assumed, and the resulting error is
not visible as an error. Silent failure of this kind is a long-established class in production
software \cite{gunawi2008eio,yuan2014simple,alquraan2018partition} and is now being catalogued for
LLM orchestration frameworks and inference engines as well \cite{xue2025agentbugs,liu2026inference}.
What changes in an evaluation setting is the artefact. The product of a silent failure is not a
degraded user session that someone reports; it is a measurement that is wrong, looks right, and is
written down.

This work reports the evaluation of one such pipeline, built for ATT\&CK mapping over static,
dynamic and network evidence, and the failures of the instrument that produced that evaluation. The
system fuses six deterministic evidence layers, three domain analysts running as ReAct loops over
tool interfaces, an LLM judge, a weighted corroboration cascade, two retrieval tiers and
deterministic confidence gating; Figure~\ref{fig:architecture} draws it. It was built around four
architectural claims: that multi-agent consensus over heterogeneous evidence channels beats a single
reader, that a multi-layer corroboration cascade improves evidence quality, that two-tier
attribution supplies useful priors, and that falsification before confidence disciplines unsupported
claims. Each claim was measured at an equal total token budget against a simpler alternative, with a
stochastic-noise control, with the firing rate of each mechanism measured before its effect, and
with intervals resampled at the level the observations are independent at.

One of those four survives in part and three do not. Decomposing the analysis into several model
calls is worth \fact{cape-negotiated-delta} F1 over a single judge, but the negotiation the design
was built around accounts for \fact{cape-mechanism-delta} of it, so what pays is the calls rather
than the argument between them. The cascade, the attribution tier and the confidence gate each
return a negative or a near-null result. Alongside them the deterministic technique assigner, the
one component measured against external ground truth rather than our own cohort, gives the cleanest
identifier gate of its three candidate backends on both corpora it was tested on.

The study that established all of this was itself wrong several times over. Seven defects in the
instrument produced plausible results rather than errors, and a suite of \fact{test-count} passing
tests caught none of them. The largest is this paper's own statistics: every interval in an earlier
draft resampled rows when the observations were clustered, and correcting that widened the anchor
interval \fact{baseline-f1-widening}-fold and turned two claims of equivalence into bounds.
